\begin{landscape}
\section{Introducción}
La sustentabilidad requiere mantener la producción sin trasladar sus costos al agua, al clima o a otras personas. Se comparan prácticas específicas, no países enteros. Los casos proceden de fuentes de distintos años: sirven para contrastar mecanismos y decisiones, no para construir una clasificación estadística de desempeño en 2026.

\section{Comparación de patrones de producción y consumo}
\subsection{Agua de riego: extracción y reutilización}
\begin{tabularx}{\linewidth}{|Y|Y|}
\hline
\rowcolor{tableblue}\textbf{MÉXICO} & \textbf{OTRO PAÍS: ISRAEL}\\
\hline
\textbf{Práctica 1. Bombeo de agua subterránea para riego sin ajustarlo a la disponibilidad.}
\par La revisión ambiental de la OCDE documentó que subsidiar la electricidad para bombeo reducía el incentivo a ahorrar agua y contribuía a la sobreexplotación de acuíferos. También identificó avances en modernización, por lo que no sería correcto afirmar que toda la agricultura mexicana utiliza riego ineficiente. Este diagnóstico corresponde a 2013 y sus antecedentes; no describe por sí solo la situación de cada acuífero actual \parencite[pp. 28, 71]{ocdeMexico2013}.
\par La práctica es insustentable cuando la extracción compromete la renovación del recurso: disminuye la reserva, eleva el costo energético del bombeo y afecta la disponibilidad para otros usuarios y ecosistemas. Instalar goteo no basta si el ahorro se usa para ampliar la superficie y aumentar el volumen total extraído.
&
\textbf{Práctica 1. Reutilización de aguas residuales tratadas en agricultura.}
\par Israel ha invertido en tratamiento y reutilización a gran escala; la OCDE lo identifica como el mayor usuario de efluentes reciclados para agricultura entre sus miembros. Sustituir parte del agua dulce por efluente tratado reduce la presión sobre fuentes convencionales cuando existen control de calidad, asignación del agua y supervisión sanitaria \parencite{ocdeIsrael2023}.
\par No es una solución sin costos: la misma revisión advierte contaminación de aguas superficiales y subterráneas y recomienda reducir nutrientes agrícolas y mejorar infraestructura. Por tanto, el referente sustentable es el reúso seguro bajo control, no cualquier aplicación de agua residual. Para México, su adaptación requiere adecuarlo al cultivo, suelo, calidad del efluente y capacidad técnica de cada comunidad.
\\\hline
\multicolumn{2}{|p{\dimexpr\linewidth-2\tabcolsep-2\arrayrulewidth\relax}|}{\textbf{ODS 6: metas 6.3 y 6.4.} El indicador \textbf{6.3.1} mide la proporción de flujos de aguas residuales domésticas e industriales tratados de manera segura; no mide directamente el porcentaje reutilizado en agricultura. El \textbf{6.4.1} sigue el cambio de eficiencia en el uso del agua y el \textbf{6.4.2} relaciona la extracción de agua dulce con los recursos disponibles, considerando los requerimientos ambientales \parencite{onuIndicadores}. Para evaluar una unidad de riego se necesitan registros de extracción, agua reutilizada, calidad y producción; estos son insumos de evaluación local, no valores nacionales oficiales de los ODS. La mejora debe combinar mayor eficiencia con menor presión total, no solo más producción por hectárea.}\\
\hline
\end{tabularx}

\clearpage
\subsection{Electricidad: dependencia fósil y transición renovable}
\begin{tabularx}{\linewidth}{|Y|Y|}
\hline
\rowcolor{tableblue}\textbf{MÉXICO} & \textbf{OTRO PAÍS: ALEMANIA}\\
\hline
\textbf{Práctica 2. Sostener el suministro eléctrico mediante combustibles fósiles sin sustituirlos progresivamente.}
\par La evaluación energética de México de 2023 describe su generación por combustibles y el papel del gas natural, junto con fuentes renovables y energía nuclear. La OCDE señaló en 2024 una participación renovable todavía baja y obstáculos regulatorios y de inversión en redes \parencite{eiaMexico2023,ocdeMexico2024}.
\par La combustión genera emisiones y mantiene dependencia del abastecimiento de combustibles; en el caso agropecuario, esto afecta bombeo, refrigeración y transformación de alimentos. La práctica insustentable no es disponer de electricidad confiable, sino perpetuar emisiones evitables y consumo ineficiente sin una transición planificada. México sí cuenta con renovables: la comparación no supone su ausencia ni atribuye todos los impactos a una sola tecnología.
&
\textbf{Práctica 2. Expandir electricidad eólica y solar con infraestructura y gestión del sistema.}
\par La agencia ambiental alemana informa que las renovables cubrieron \textbf{55.1\% del consumo eléctrico bruto en 2025}; la energía eólica fue la principal fuente de su mezcla eléctrica. Esta expansión muestra que la sustitución de generación fósil puede tener una escala significativa \parencite{uba2026}.
\par La mejora exige redes, flexibilidad de demanda, almacenamiento y respaldo. No implica que Alemania carezca de combustibles fósiles ni que la energía renovable tenga impacto ambiental nulo. En México, la aplicación debe combinar proyectos bien ubicados con consulta social, acceso a financiamiento y eficiencia de equipos. Para una explotación agrícola, el bombeo solar tampoco justifica extraer más agua de la disponible: una solución energética puede agravar otro problema si no se evalúa conjuntamente.
\\\hline
\multicolumn{2}{|p{\dimexpr\linewidth-2\tabcolsep-2\arrayrulewidth\relax}|}{\textbf{ODS 7: metas 7.2 y 7.3.} El indicador \textbf{7.2.1} es la participación de las renovables en el \emph{consumo final total de energía}; incluye más usos que la electricidad. Por eso, el 55.1\% eléctrico alemán NO es un valor de 7.2.1. La UBA reporta por separado 23.8\% de consumo final bruto renovable bajo la metodología europea RED, que tampoco debe equipararse automáticamente al indicador ONU. El \textbf{7.3.1} mide intensidad energética en términos de energía primaria y PIB \parencite{onuIndicadores,uba2026}. En una unidad agrícola, kWh por tonelada y consumo por operación permiten seguimiento operativo, pero no son sustitutos de esos indicadores nacionales.}\\
\hline
\end{tabularx}

\clearpage
\subsection{Envases: descarte mezclado y retorno organizado}
\begin{tabularx}{\linewidth}{|Y|Y|}
\hline
\rowcolor{tableblue}\textbf{MÉXICO} & \textbf{OTRO PAÍS: ALEMANIA}\\
\hline
\textbf{Práctica 3. Consumir envases de un solo uso y desecharlos mezclados con otros residuos.}
\par Mezclar envases recuperables con restos orgánicos dificulta su separación y aprovechamiento. El diagnóstico de SEMARNAT de 2020 explica la importancia de la recolección selectiva y documenta limitaciones en infraestructura de tratamiento; además, reconoce la falta de cifras oficiales completas sobre la recuperación mediante pepena \parencite[secs. 2.2.1.3 y 2.2.3.2]{semarnat2020}.
\par No se infiere que México no recicle ni que todos los materiales tengan la misma tasa. El problema es el patrón de usar y descartar sin separación y sin destino comprobable. La pérdida de material aprovechable aumenta la demanda de materias primas y la carga de disposición final. Los envases de agroquímicos requieren manejo específico y no deben mezclarse con los de bebidas ni incluirse sin más en un sistema doméstico de retorno.
&
\textbf{Práctica 3. Depósito y devolución de envases de bebidas.}
\par El sistema alemán DPG articula productores, importadores, comercios y operadores. Los envases sujetos al sistema llevan identificación; el consumidor paga un depósito de \textbf{0.25 euros} y lo recupera al devolverlos. Los materiales recogidos se clasifican y se destinan al reciclaje mediante una cadena organizada \parencite{dpgProceso}.
\par La devolución da valor al retorno y favorece flujos más separados. Sin embargo, devolver no equivale a reutilizar la misma botella, ni una tasa de retorno equivale a la tasa nacional de reciclaje. Esta experiencia no cubre todos los residuos. Adaptarla en México requeriría obligaciones claras, puntos accesibles, trazabilidad y participación remunerada de recuperadores; prevenir envases innecesarios y favorecer opciones reutilizables debe preceder al reciclaje.
\\\hline
\multicolumn{2}{|p{\dimexpr\linewidth-2\tabcolsep-2\arrayrulewidth\relax}|}{\textbf{ODS 12: meta 12.5.} El indicador \textbf{12.5.1} mide la tasa nacional de reciclaje y las toneladas de material reciclado \parencite{onuIndicadores}. Para evaluar un sistema local de envases se deben distinguir unidades vendidas, unidades devueltas, masa recuperada y masa efectivamente reciclada. Recolección, valorización energética, reúso y reciclaje no son sinónimos. La meta también exige prevención y reducción, por lo que aumentar toneladas recicladas sin reducir residuos no prueba por sí solo mayor sustentabilidad.}\\
\hline
\end{tabularx}
\end{landscape}

\section{Conclusión}
El comparativo muestra tres mecanismos distintos: reducir la presión sobre el agua mediante asignación y reúso seguro; sustituir generación fósil con renovables y eficiencia; y cambiar el descarte de envases mediante incentivos, separación y trazabilidad. Su valor no reside en copiar un país, sino en identificar qué condiciones hacen funcionar cada práctica y qué riesgos permanecen. Las advertencias de la OCDE sobre contaminación en Israel y la distinción entre electricidad y energía total en Alemania evitan presentar esas experiencias como modelos perfectos.

Desde la perspectiva de la ingeniería agronómica, considero prioritario vincular la decisión tecnológica con límites ambientales y viabilidad social. Un sistema de riego eficiente necesita medición de extracción y criterios de calidad; un equipo eléctrico eficiente necesita operar dentro de un plan de uso de agua; y un esquema de retorno requiere destinos materiales comprobables y condiciones dignas para quienes recuperan residuos. No bastaría con instalar infraestructura si faltan mantenimiento, capacitación, vigilancia o acceso para pequeños productores.

En consecuencia, propondría comenzar con diagnósticos locales, metas medibles y pilotos evaluados antes de ampliar las inversiones. Los indicadores ODS ofrecen un marco para esa evaluación, pero sus denominadores y escalas deben conservarse. Mi conclusión es que la transición mexicana debe combinar prevención, eficiencia y responsabilidad compartida, con decisiones ajustadas al territorio y evidencia verificable, en lugar de promesas tecnológicas aisladas.\footnote{Se utilizó AulaTeX con un modelo de lenguaje para apoyar la organización, redacción y revisión del documento. La versión se contrastó con las fuentes citadas; corresponde al estudiante revisar las afirmaciones y asumir o ajustar la reflexión personal antes de entregar.}
